%% kapa: exam=1 round=32 year=2021 period=1 form=A subject=real_estate subject_ko=부동산학원론 seq=28 number=108
%% source: Q-Net 감정평가사 기출문제 (cst00309) artlSeq=5208397 file=extracted/1008/5208397/A형/2021년도 제32회 감정평가사 1차 1교시 A형.pdf
%% conversion: kordoc 4.24.51 -> tools/gen1.py (auto); review=proofread-sonnet; flags=none
\begin{question}{108}
부동산개발사업에 관련된 설명으로 옳은 것을 모두 고른 것은?
\begin{framedtext}
ㄱ. 개발기간의 연장, 이자율의 인상, 인플레이션의 영향으로 개발비용이 증가하는 위험은 비용위험에 속한다. \par
ㄴ. 개발부동산의 선분양제도는 후분양제도에 비해 사업시행자가 부담하는 시장위험을 줄일 수 있다. \par
ㄷ. 민감도분석에 있어 주요 변수로는 토지구입비, 개발기간, 분양가격 등이 있다. \par
ㄹ. 수익성지수가 1보다 크다는 것은 순현가가 ‘0’(zero)보다 크다는 뜻이다. \par
\end{framedtext}
\begin{choices}
  \item ㄱ, ㄴ
  \item ㄴ, ㄷ
  \item ㄱ, ㄷ, ㄹ
  \item ㄴ, ㄷ, ㄹ
  \item ㄱ, ㄴ, ㄷ, ㄹ
\end{choices}
\end{question}
\answer{5}
